%========================================================================
% Modelo para elaboracao de textos academicos: TCC, dissertacoes e teses
% Elaborado pelo GISIS - Grupo de Imageamento Sismico e Inversão Sismica.
%========================================================================
\chapter{Revisão Bibliográfica}
\label{ch:revisaobibliografica}

\section{Modelagem sísmica}




\subsection*{Equação da onda para meios acústicos}




\subsection*{Equação analítica para ondas refratadas}




\subsection*{Equação \textit{eikonal}}




\subsection*{Método clássico}




\subsection*{\textit{Fast Iterative Method}}




\subsection*{\textit{Fast Sweeping Method}}




\section{Inversão tomográfica}

% Tipos de tomografia (reflexão, difração, transmissão e refração)





\subsection*{Tomografia de refração}
% Discretização do modelo
% Raios ilustrativos
% Função objetivo e linearização
% Regularização 
% Resolução do sistema linear iterativamente

\subsection*{Obtenção do dado observado}
% picking manual
% picking por aprendizado de máquina
% picking por equações algébricas que relacionam energia (sta/lta - espectro) 

 
 
 
 
 
 
 
 
 
 
 
 
%Escreva aqui sua revisão de literatura. 
%
%\section{Exemplo de seção}
%\label{sec:exemplo}
%
%Na seção \ref{sec:exemplo}
%
%\subsection{Exemplo de subseção}
%
%Ao comentar sobre sua figura no texto, é importante usar o comando \verb_\ref{LabelParaCitarNoTexto}_. Ele também é utilizado para fazer referência a Tabelas, Equações, Capítulos, etc. e vai atribuir a numeração adequada ao elemento que possui aquele \textit{label}. Um exemplo de texto mencionando a figura é dado a seguir.
%
%Na Figura \ref{LabelParaCitarNoTexto}, pode-se observar um quadrado azul utilizado para exemplificar onde vai entrar a imagem.
%\begin{figure}[!h]
%    \centering
%    \includegraphics[width=.65\textwidth]{Imgs/RevisaoBibliografica/ExemploImagem.png}
%    \caption[Titulo da imagem, como deve aparecer na lista.]{Na legenda da imagem, deve-se descrever exatamente o que se observa nela. A legenda pode se tornar um pouco longa, mas é importante incluir todas as informações que forem necessárias para seu entendimento. Lembre-se de incluir a fonte de onde a imagem foi selecionada ou adaptada/modificada. Se for uma imagem própria, incluir a fonte como: ``Fonte: O autor".}
%    \label{LabelParaCitarNoTexto}
%\end{figure}
%
%Também é possível colocar mais de uma imagem numa mesma figura, como na Figura \ref{duasfiguras}. Você pode citá-las individualmente utilizando o \textit{label} atribuído a cada uma delas. Na Figura \ref{primeirasubfig}, temos a primeira imagem e, na Figura \ref{segundasubfig}, temos a segunda imagem.
%
%\begin{figure}[h]
%	\centering
%	\subfloat[]{
%	   \label{primeirasubfig}
%		\includegraphics[width= 0.45\textwidth]{Imgs/RevisaoBibliografica/ExemploImagem.png}
%	}
%	\subfloat[]{
%	   \label{segundasubfig}
%		\includegraphics[width= 0.45\textwidth]{Imgs/RevisaoBibliografica/ExemploImagem.png}
%	}
%	%O que esta entre [] e o titulo da imagem da maneira como deve a aparecer na lista de ilustracoes. Entre {}, e a legenda completa da Figura.
%	\caption [Exemplo de Figura com mais de uma imagem.]{Exemplo de Figura com mais de uma imagem. a) Primeira imagem. b) Segunda imagem. Fonte: o autor.}
%	\label{duasfiguras}
%\end{figure}
%
%%Exemplo de Equação
%A famosa fórmula de Bhaskara é dada pela Equação \ref{eq:bhaskara}.
%
%\begin{equation}
%    x = \frac{-b \pm \sqrt{b^{2}-4ac}}{2a}
%    \label{eq:bhaskara}
%\end{equation}
